\documentclass[12pt,a4paper]{report}

\usepackage[utf8]{inputenc}
\usepackage[T1]{fontenc}
\usepackage[english]{babel}
\usepackage[a4paper,margin=3cm]{geometry}
\usepackage{setspace}
\usepackage{parskip}
\usepackage{titlesec}
\usepackage{enumitem}
\usepackage{amsmath}
\usepackage{amssymb}
\usepackage{hyperref}
\usepackage{listings}
\usepackage{xcolor}

% Line spacing
\setstretch{1.25}

% Chapter format — numbered, large, bold, no "Chapter X" prefix
\titleformat{\chapter}[block]
  {\normalfont\LARGE\bfseries}
  {\thechapter.}{0.5em}{}

\titlespacing*{\chapter}{0pt}{0pt}{1.2cm}

% Section format
\titleformat{\section}
  {\normalfont\large\bfseries}
  {\thesection}{0.5em}{}

% Code listing style
\lstset{
  language=Python,
  basicstyle=\ttfamily\small,
  keywordstyle=\color{blue},
  commentstyle=\color{gray},
  stringstyle=\color{orange},
  numbers=left,
  numberstyle=\tiny\color{gray},
  breaklines=true,
  frame=single,
  rulecolor=\color{black!20},
  backgroundcolor=\color{black!3},
  tabsize=4
}

% Hyperlink styling
\hypersetup{
  colorlinks=true,
  linkcolor=black,
  urlcolor=blue
}

\begin{document}

% ─── Title Page ───────────────────────────────────────────────────────────────
\begin{titlepage}
    \centering

    \vspace*{3cm}

    {\Huge\bfseries VHGRAD \par}

    \vspace{0.8cm}

    {\Large Development of a Didactic Deep Learning Framework from Scratch \par}

    \vspace{3cm}

    {\large \textbf{Virvi Huta} \par}

    \vfill

    {\large 2026 \par}
\end{titlepage}

\clearpage

% ─── Table of Contents ────────────────────────────────────────────────────────
\tableofcontents
\clearpage

% ─── Technical Progression ────────────────────────────────────────────────────
\chapter*{Technical Progression}
\addcontentsline{toc}{chapter}{Technical Progression}

\begingroup
\small

\begin{enumerate}[
    leftmargin=0pt,
    itemindent=1.5em,
    label=\textbf{\arabic*.},
    itemsep=0.75em,
    topsep=0.3em,
    parsep=0pt
]

\item \textbf{Foundations of Neural Computation}\\
Fundamental concepts underlying modern neural networks.

\item \textbf{Building the First Components}\\
Single neurons, multi-neuron layers, tensors and matrix operations with NumPy.

\item \textbf{Structured Layer Design}\\
Encapsulation of dense layers into reusable classes with organized data handling.

\item \textbf{Non-Linearity through Activation}\\
Step, linear, sigmoid, ReLU and Softmax activation functions and their roles.

\item \textbf{Measuring Network Error}\\
Cross-entropy loss formulation and accuracy metrics for classification.

\item \textbf{Parameter Optimization}\\
First principles of gradient-based learning and iterative parameter updates.

\item \textbf{Calculus of Learning}\\
Numerical and analytical derivatives, partial derivatives and the chain rule.

\item \textbf{Gradient Flow through the Network}\\
Full backward pass with analytically derived gradients across layers and losses.

\item \textbf{Advanced Optimization Methods}\\
SGD, learning rate scheduling, momentum and adaptive optimizers including Adam.

\item \textbf{Controlling Generalization}\\
Weight regularization, dropout and validation strategies to reduce overfitting.

\item \textbf{Extending to Other Task Types}\\
Binary classification and regression with appropriate outputs and loss functions.

\item \textbf{Training on Real Data}\\
End-to-end pipeline covering loading, preprocessing, batching and model training.

\item \textbf{Deployment and Reuse}\\
Parameter serialization, model loading and inference on unseen data.

\end{enumerate}

\endgroup

\clearpage

\chapter{Foundations of Neural Computation}

\section{The Single Neuron}

A neuron is the fundamental computational unit of a neural network. It takes a set of inputs,
multiplies each by a corresponding weight, sums the results and adds a scalar offset called a bias.
Given inputs $x_1, x_2, \ldots, x_n$, weights $w_1, w_2, \ldots, w_n$ and bias $b$,
the output of a single neuron is:

\begin{equation}
    z = \sum_{i=1}^{n} x_i w_i + b
\end{equation}

The weights control how strongly each input influences the output. A large positive weight
amplifies the corresponding input, a negative weight inverts it and a weight near zero makes
that input nearly irrelevant. The bias shifts the output independently of the inputs.

The first implementation computed this explicitly for four inputs with no abstraction:

\begin{lstlisting}[language=Python]
inputs  = [1, 2, 3, 2.5]
weights = [0.2, 0.8, -0.5, 1.0]
bias    = 2

output = (inputs[0] * weights[0] +
          inputs[1] * weights[1] +
          inputs[2] * weights[2] +
          inputs[3] * weights[3] +
          bias)
# output: 4.8
\end{lstlisting}

Substituting the values directly:
$1 \times 0.2 + 2 \times 0.8 + 3 \times (-0.5) + 2.5 \times 1.0 + 2 = 4.8$

Writing every multiplication by hand was intentional — the goal was to see the raw mechanics
before introducing any abstraction.

% ─────────────────────────────────────────────────────────────────────────────
\section{A Layer of Neurons}

A layer is a collection of neurons that all receive the same input vector but each apply
their own weights and bias. For a layer of $m$ neurons with shared input
$\mathbf{x} \in \mathbb{R}^n$, the output of neuron $j$ is:

\begin{equation}
    z_j = \sum_{i=1}^{n} x_i w_{ji} + b_j, \quad j = 1, \ldots, m
\end{equation}

The implementation first wrote all three outputs explicitly — every multiplication typed out —
to make the repetitive structure visible. Then it was refactored into a loop:

\begin{lstlisting}[language=Python]
inputs  = [1, 2, 3, 2.5]
weights = [[0.2, 0.8, -0.5, 1.0],
           [0.5, -0.91, 0.26, -0.5],
           [-0.26, -0.27, 0.17, 0.87]]
biases  = [2, 3, 0.5]

layer_outputs = []
for neuron_weights, neuron_bias in zip(weights, biases):
    neuron_output = 0
    for n_input, weight in zip(inputs, neuron_weights):
        neuron_output += n_input * weight
    neuron_output += neuron_bias
    layer_outputs.append(neuron_output)

# layer_outputs: [4.8, 1.21, 2.385]
\end{lstlisting}

The loop makes clear that a layer is nothing more than the same neuron computation
repeated across rows of the weight matrix. The structure — same inputs, different weights
per neuron — is the core of what a dense layer is.

% ─────────────────────────────────────────────────────────────────────────────
\section{The Dot Product}

Before moving to NumPy it was important to recognise that a neuron's weighted sum is
mathematically a dot product. Given two vectors $\mathbf{a}, \mathbf{b} \in \mathbb{R}^n$,
their dot product is:

\begin{equation}
    \mathbf{a} \cdot \mathbf{b} = \sum_{i=1}^{n} a_i b_i
\end{equation}

This was verified manually and then with a loop for $\mathbf{a} = [1, 2, 3]$
and $\mathbf{b} = [2, 3, 4]$:

\begin{equation}
    1 \times 2 + 2 \times 3 + 3 \times 4 = 20
\end{equation}

Recognising this equivalence is what makes the transition to NumPy clean —
a neuron's computation is exactly one call to a dot product function.

% ─────────────────────────────────────────────────────────────────────────────
\section{Vectorized Computation with NumPy}

Writing loops in Python is slow. NumPy replaces them with highly optimised C routines
that operate on entire arrays at once. The single neuron collapses to one line:

\begin{lstlisting}[language=Python]
outputs = np.dot(inputs, weights) + bias
# outputs: 4.8
\end{lstlisting}

For the full layer the weight matrix is passed first so NumPy computes one dot product
per row — one output per neuron:

\begin{lstlisting}[language=Python]
outputs = np.dot(weights, inputs) + biases
# outputs: [4.8, 1.21, 2.385]
\end{lstlisting}

The argument order matters. \texttt{np.dot(weights, inputs)} treats each row of the
weight matrix as one neuron's weight vector and computes its dot product with the input
vector, producing one scalar per neuron.

% ─────────────────────────────────────────────────────────────────────────────
\section{Matrices and Shape Alignment}

A plain Python list passed to NumPy produces a 1D array — a vector with no orientation.
Wrapping it in an extra list creates a 2D row vector of shape $(1, n)$, and transposing
it with \texttt{.T} produces a column vector of shape $(n, 1)$. This was verified
with a matrix dot product:

\begin{lstlisting}[language=Python]
a = np.array([[1, 2, 3]])     # shape (1, 3)
b = np.array([[2, 3, 4]]).T   # shape (3, 1)
np.dot(a, b)                  # [[20]], shape (1, 1)
\end{lstlisting}

The result is a $(1, 1)$ matrix rather than a scalar because matrix multiplication
was performed rather than a vector dot product. Understanding how shapes interact
under multiplication is essential for everything that follows.

% ─────────────────────────────────────────────────────────────────────────────
\section{Batched Forward Pass}

Processing one sample at a time is inefficient. In practice a neural network processes
a batch of samples simultaneously. If the input matrix $X$ has shape
$(N, n)$ — $N$ samples each with $n$ features — and the weight matrix $W$ has shape
$(m, n)$ — $m$ neurons each with $n$ weights — then the full layer output is:

\begin{equation}
    Z = X W^T + \mathbf{b}
\end{equation}

where $\mathbf{b} \in \mathbb{R}^m$ is the bias vector broadcast across all samples,
and $Z$ has shape $(N, m)$ — one output per neuron per sample.

The transpose $W^T$ flips $W$ from $(m, n)$ to $(n, m)$ so that the matrix product
$X W^T$ of shapes $(N, n) \cdot (n, m)$ is defined and yields $(N, m)$.

\begin{lstlisting}[language=Python]
inputs  = [[1.0, 2.0, 3.0, 2.5],
           [2.0, 5.0, -1.0, 2.0],
           [-1.5, 2.7, 3.3, -0.8]]

weights = [[0.2, 0.8, -0.5, 1.0],
           [0.5, -0.91, 0.26, -0.5],
           [-0.26, -0.27, 0.17, 0.87]]

biases = [2.0, 3.0, 0.5]

layer_outputs = np.dot(inputs, np.array(weights).T) + biases
\end{lstlisting}

The result is a $(3, 3)$ matrix — three samples, three neuron outputs each:

\begin{equation}
    Z = \begin{pmatrix}
        4.800 & 1.210 & 2.385 \\
        8.900 & -1.810 & 0.200 \\
        1.410 & 1.051 & 0.026
    \end{pmatrix}
\end{equation}

Each row corresponds to one input sample and each column to one neuron.
This is a complete batched forward pass through a single dense layer —
the foundation on which every subsequent component of VHgrad is built.

\end{document}